\begin{abstract}
多模态大模型赋能的空天信息处理，为灾后应急侦察提供了语言驱动、感知与规划闭环的新路径。
本文给出一套已实现的空天应急侦察智能体与地理配准仿真平台：操作员以自然语言下达任务，视觉语言大模型联合双时相损伤分类与场景分割处理跨时相空天影像，智能体在由灾前/灾后卫星影像与连续底图合成的俯视仿真中闭合感知--规划--执行回路，并提供实时人机协同控制台。

观测几何是本文核心工程贡献之一：固定传感器与固定视场角下，巡航 $1330\,\mathrm{m}$（$3\times3$ 瓦片，$1.50\,\mathrm{m/px}$）下降到 $443\,\mathrm{m}$（恰好一整张瓦片，$0.50\,\mathrm{m/px}$），在连续底图上收缩视场而不是对同一裁块撤销高斯模糊。
信息上界仍是源影像原生 $0.5\,\mathrm{m/px}$，不是厘米级低空纹理；题目锚定地理 ROI，答案作用域不随高度改变。

在此平台上，我们把``第一次俯视是否足够''表述为\textbf{预算受限的主动重观测}：以有效 GSD 为物理可审计量，把证据是否足够转成一次可审计的下降--居中动作，并用``可识别性优先''的判据检验该机制。
在新观测模型与一个已披露的强感知底座下，降高确实改变了证据：逐建筑 $n_{\mathrm{flip}}=733$、$n_{\mathrm{correctable}}=375$，但翻转有得有失，题目层面可纠正 $13$ 题、破坏 $10$ 题（MARGINAL）——机制可测、效应小且不单调，支撑``重观测何时值得''这一问题的存在性，而非``重观测改善成功率''。

感知底座采用 xView2 冠军架构的预训练权重，其训练划分见过评测事件（leaky）；本文将其明确降级为公开披露的感知组件，不再主张跨事件泛化。
一个三方对照测量了``损伤分级坍缩''的根源：同架构事件不相交重训后损伤 F1 仅 $0.017$，现成权重为 $0.630$，差距约 $97\%$ 来自事件曝光、约 $3\%$ 来自架构——即瓶颈是跨事件泛化，而非架构或训练规模。

本文贡献分两层：可运行的灾害智能体、物理自洽的仿真环境与预算受限主动重观测机制是工程与机制存在性声明；可识别性优先评估、对负结果的反事实排除与瓶颈分解是科学结论。
二者不得混读为复检已经提高任务成功率。
\end{abstract}
